\section{The Companions of Christ: St.~Stephen, St.~John, and the Holy Innocents}\label{sec:comites}

Before the candles of Christmas Day have burned down, the 1962 Missal turns from
the straw of Bethlehem to stones, exile, and spilled blood. On 26, 27, and
28~December, the Church keeps three second-class feasts directly inside the
Christmas Octave: St.~Stephen the Protomartyr, St.~John the Apostle and
Evangelist, and the Holy Innocents. Medieval liturgists called them the
\textit{Comites Christi}---the closest companions surrounding the newborn King.

Preaching on the Feast of the Holy Innocents, St.~Bernard of Clairvaux explained
why these three feasts stand together as a single continuation of Christmas:

\begin{quote}
You may have noticed that in the same three solemnities we celebrate three
distinct kinds of sanctity. Nor will it be easy to discover any human sanctity
which does not belong to one of the three. For in St.~Stephen we have both the
will and the act of martyrdom; in St.~John Evangelist we have the will without
the act; whilst in the Holy Innocents we have the act without the will. All of
them drank the chalice of salvation, either in act and will together, or in will
alone, or in act alone.
\end{quote}

In Book~VII, Chapter~42 of his \textit{Rationale}, William Durandus joined the
same three companions to Canticles~5:10, ``My beloved is white and ruddy, chosen
out of thousands'': ruddy in St.~Stephen's blood, white in St.~John's virginal
purity, and chosen out of thousands in the white-robed host of the Innocents.

Under the 1960 \textit{Codex Rubricarum}, each of these three feasts ranks as a
second-class feast (General Rubric~91, row~16) that also stands as the second,
third, and fourth day within the Christmas Octave (RG~68~a--c). On Monday
through Saturday, the saint's Mass is celebrated with the \textbf{privileged
commemoration of the Octave of the Nativity} (RG~109~c) at both Low Mass and
Sung Mass, using the Day Mass orations (\textit{Concede}, \textit{Oblata},
\textit{Praesta}), the Creed, the Nativity Preface (overriding even the Apostles
Preface on St.~John's day), and the proper \textit{Communicantes} (\textit{diem
sacratissimum}). Whenever 26, 27, or 28~December falls on a \textbf{Sunday},
the Sunday within the Octave takes precedence (row~15 over row~16), and the
saint's feast receives an ordinary commemoration at Low Mass only (omitted in a
non-conventual Sung Mass, RG~111~b).

\subsection{26 December: St.~Stephen, Protomartyr}

On 26~December the priest puts on \textbf{red} vestments, and the Roman station
moves to the round basilica of \textit{Sanctus Stephanus in Coelio monte},
St.~Stephen on the Caelian Hill. The Introit chants Psalm~118~[119]:23,
\textit{Sederunt principes, et adversum me loquebantur}, ``Princes sat, and
spoke against me: and the wicked persecuted me.'' The Collect turns adoration
into imitation:

\begin{massbox}{Collect of St.~Stephen (26 December)}
\textbf{Latin} (1962 Roman Missal). Da nobis, quaesumus, Domine, imitari quod
colimus: ut discamus et inimicos diligere; quia ejus natalicia celebramus, qui
novit etiam pro persecutoribus exorare Dominum nostrum Jesum Christum, Filium
tuum: Qui tecum vivit et regnat in unitate.

\textbf{English} (Lefebvre, 1925). Grant us, we beseech Thee, O Lord, so to
imitate what we revere that we may learn to love even our enemies: for we
celebrate the day of his birth to immortality, who could even plead with our
Lord Jesus Christ Thy Son on behalf of his persecutors. Who with Thee.
\end{massbox}

The Lesson (Acts~6:8--10; 7:54--59) recounts Stephen full of grace and
fortitude, gazing up into the opened heavens to see Jesus standing at the right
hand of God, crying ``Lord Jesus, receive my spirit,'' and falling on his knees
to pray ``Lord, lay not this sin to their charge'' before falling asleep in the
Lord. The Gospel (Matt~23:34--39) records Christ's own lament over Jerusalem,
which kills the prophets and stones those sent to her, yet longs to gather her
children as a hen gathers her chicks under her wings.

In the Second Nocturn of St.~Stephen's Matins, St.~Fulgentius of Ruspe sets
25~December and 26~December face to face in a sentence that medieval liturgists
never tired of repeating: ``Yesterday our King, clothed in the robe of flesh,
came forth from the palace of the Virgin's womb to visit the world; to-day His
soldier, quitting the tabernacle of the body, enters triumphant into heaven.''
St.~Peter Chrysologus observes in Sermon~154 that just as Peter took his name
from the rock (\textit{petra}) on which the Church stands firm, so Stephen took
his name from the crown (\textit{stephanos}), because he was the first to win
the crown of martyrdom for the name of Christ. Preaching on Acts~7 in Homily~18,
St.~John Chrysostom watches Stephen mirror every gesture of Calvary---cast
outside the city, commending his spirit, and forgiving his murderers---so that
``even in death, Confession and Preaching'' shine forth without a trace of
bitterness.

\subsection{27 December: St.~John, Apostle and Evangelist}

On 27~December the vestments return to \textbf{white}, and the Roman station
returns to \textit{Sancta Maria maior}, uniting the Virgin Mother with the
beloved disciple to whom Jesus entrusted her from the cross. The Introit and
Lesson come from Ecclesiasticus~15:1--6: \textit{In medio Ecclesiae aperuit os
ejus}, ``In the midst of the church the Lord shall open his mouth, and shall
fill him with the spirit of wisdom and understanding, and shall clothe him with
a robe of glory.''

\begin{massbox}{Collect of St.~John the Evangelist (27 December)}
\textbf{Latin} (1962 Roman Missal). Ecclesiam tuam, Domine, benignus illustra:
ut, beati Joannis Apostoli tui et Evangelistae illuminata doctrinis, ad dona
perveniat sempiterna. Per Dominum nostrum.

\textbf{English} (Lasance, 1945). Of Thy loving-kindness, O Lord, shed light
upon Thy Church; that, being enlightened by the teachings of blessed John, Thine
apostle and evangelist, it may attain to Thine everlasting gifts. Through our
Lord.
\end{massbox}

The Gospel (John~21:19--24) stands at the lakeside after the Resurrection:
Christ foretells Peter's crucifixion (``Follow me''), and when Peter looks back
at the disciple who leaned on Jesus's breast at supper and asks ``Lord, and what
shall this man do?'' the Lord answers, ``So I will have him to remain till I
come, what is it to thee? Follow thou me.'' In his 124th Tractate on John
(section~5, read at St.~John's Matins), St.~Augustine sees in Peter and John the
two lives known to the Church: one in faith, temporal labor, and active service
(signified by Peter following Christ through the Passion), and the other in
sight, repose, and the wages of contemplation (signified by John remaining until
the Lord comes). Every Christian shares in both: active fidelity now, looking
toward the vision that abides forever.

At Matins, the Sixth Lesson adds St.~Jerome's unforgettable portrait from his
\textit{Commentary on Galatians} of the aged Apostle in Ephesus:

\begin{quote}
The Blessed Evangelist John lived at Ephesus down to an extreme old age, and, at
length, when he was with difficulty carried to the Church, and was not able to
exhort the congregation at length, he was used simply to say at each meeting,
``My little children, love one another.'' At last the disciples and brethren
were weary with hearing these words continually, and asked him, ``Master,
wherefore ever sayest thou this only?'' Whereto he replied to them, (worthy of
John,) ``It is the commandment of the Lord, and if this only be done, it is
enough.''
\end{quote}

Because 27~December lies inside the Christmas Octave, the Nativity Preface
overrides the Apostles Preface---reminding the assembly that the eagle who
soared highest into ``In the beginning was the Word'' rested his head on the
human breast of the Word made flesh.

\subsection{28 December: The Holy Innocents, Martyrs}

On 28~December, the Roman station moves to \textit{Sanctus Paulus}, St.~Paul
Outside the Walls. Here readers familiar with older pre-1960 hand missals
encounter one of the clearest liturgical changes of the 1962 Roman Rite. In
older books, when 28~December fell on a weekday, the Church vested in violet of
mourning, omitted the \textit{Gloria}, \textit{Te Deum}, and \textit{Alleluia}
(substituting a Tract), wore folded chasubles, and closed Mass with
\textit{Benedicamus Domino}, reserving red and joyful chants for Sundays. Under
the 1960 \textit{Codex Rubricarum} and the 1962 Missal, every one of those
weekday mourning exceptions was abolished: folded chasubles disappeared
(RG~137), \textbf{red} vestments are worn unconditionally (RG~124~e), the
\textit{Gloria} and \textit{Te Deum} are sung, the annual Mass sings the
\textbf{Alleluia} (\textit{Laudate, pueri, Dominum}, Ps~112~[113]:1), and Mass
ends with \textit{Ite, missa est} (RGMR~507). By a quiet irony of precedence,
when 28~December falls on a Sunday in 1962, the Sunday within the Octave wins
and the Holy Innocents are only commemorated in Low Mass---so whenever their own
Mass is celebrated on Monday through Saturday, it is now vested in the red of
martyr triumph.

The Introit opens with Psalm~8:3, \textit{Ex ore infantium, Deus, et lactentium
perfecisti laudem}, ``Out of the mouth of infants and of sucklings thou hast
perfected praise.'' The Collect faces the scandal of infant suffering directly:

\begin{massbox}{Collect of the Holy Innocents (28 December)}
\textbf{Latin} (1962 Roman Missal). Deus, cujus hodierna die praeconium
Innocentes Martyres non loquendo, sed moriendo confessi sunt: omnia in nobis
vitiorum mala mortifica; ut fidem tuam, quam lingua nostra loquitur, etiam
moribus vita fateatur. Per Dominum nostrum.

\textbf{English} (Lasance, 1945). O God, Whose praise the martyred innocents did
this day proclaim, not by speaking, but by dying, do to death in us all the
malice of sinfulness, that our lives may also proclaim Thy faith, which our
tongues profess. Through our Lord.
\end{massbox}

The Lesson (Apoc~14:1--5) lifts the slaughtered boys of Bethlehem straight to
Mount Sion, singing the new canticle before the Lamb as ``the firstfruits to God
and to the Lamb,'' in whose mouth no lie was found. The Gospel (Matt~2:13--18)
brings us back to earth: Joseph warned in sleep to flee by night into Egypt with
the Child and his Mother, Herod's fury slaying every male child in Bethlehem
two years old and under, and Jeremiah's cry echoed in the Communion chant: ``A
voice in Rama was heard, lamentation and great mourning; Rachel bewailing her
children, and would not be comforted, because they are not.''

Why did the Almighty Child flee into Egypt instead of striking Herod down? In
his eighth Homily on Matthew, St.~John Chrysostom answers that had Christ
acted constantly with naked divine power, men would have denied the reality of
his human flesh---the flight into Egypt proves that the Word truly entered our
vulnerability; and in Homily~9, Chrysostom makes plain that ``Christ was not
the cause of their slaughter, but the king's cruelty.'' In the Matins sermon
attributed to St.~Augustine, the Church blesses Bethlehem as ``worthy to offer
at once to God a whole white-robed army of guileless martyrs,'' calling them
``the flowers of martyrdom; hardly had these buds of the Church shown their
heads above the soil, in the winter of unbelief, when the frost of persecution
nipped them.'' At Lauds, the fourth-century Spanish poet Prudentius puts that
image into the immortal hymn \textit{Salvete flores Martyrum}
(Appendix~\ref{app:prayers}):

\begin{quote}
You, tender flock of lambs, we sing,\\
First victims slain for Christ your King:\\
Beneath the altar's heavenly ray,\\
With martyr palms and crowns ye play.
\end{quote}
